
\section{Algorithm Proposal}
\label{sec:algorithm}

As shown in the previous section, the position of a proposer does impact consensus performance. From this observation we devise an approach to dynamically adjust the frequency that each participant becomes proposer, favoring faster ones, and thus enhancing overall consensus performance. The proposed approach poses the following assumptions to the underlying consensus algorithm:  
\begin{description}
    \item[A.1] it uses rotating proposer and provides means to record the future sequence of proposers (or parameters to calculate it); 
    \item[A.2] it has a deterministic function used at each node to calculate the next proposer.   We call this function \textit{nextProposer}.  For instance, with PoS, the next proposer is a function of the known stakes from each participant. 
\end{description}

The proposed mechanisms should ensure the following basic properties:
\begin{description} 

    \item[P.1] the underlying consensus algorithm properties are preserved;     
    %For instance, it does not modify the composition of the consensus committee, unlike some of the related works; 

    \item[P.2] disregarding stake, or with a uniform stake distribution, proposer nodes with faster (slower) decisions become proposers proportionally more (less) often;

    \item[P.3] changes in a node's relative performance over time lead to respective adaptation of its frequency as proposer;  

    \item[P.4] every node eventually becomes proposer.
\end{description}

And the byzantine fault-tolerance properties:
\begin{description} 
   
    \item[P.5] dishonest nodes do not impose an arbitrary sequence of proposers;
    
    \item[P.6] dishonest nodes do not hinder the differentiation of faster/slower nodes. 

\end{description}

To achieve the desired properties, three concurrent, continuous mechanisms are proposed:
\begin{description} 
    \item[M.1] every node samples the latency and locally classifies each other node;
    
    \item[M.2] every node collaborates to generate consistent global classifications;
    
    \item[M.3] every node uses the global classification to calculate and adopt a new sequence of proposers.   
\end{description}

\subsection{M.1 - Sampling and classifying nodes}

While the techniques for sampling and classifying nodes may vary, we describe here one with low computational cost, yet effective in our experiments. 
%
Each node continuously samples the latency of the successive consensus heights.  This measure is local, comprising the time elapsed from the start of the current consensus (e.g. Tendermint's propose) to it's decision (e.g. Tendermint's decided).   Since each consensus height is associated to a known proposer, the sample denotes that proposer's performance.  

\paragraph{Set of latency samples $L$.} Let $L$ be a set with the $k$ last consensus latency samples taken by a node, $k$ being a configuration parameter.  At any time, the average consensus latency $\bar{L}$ is known at the node.

\paragraph{A parameter $T \in {0..1}$} defines a threshold for node classification against the average latency $\bar{L}$. A latency sample $s$ for a consensus instance can be considered $fast$ if $s < \bar{L} \times (1-T)$; $slow$ if $s > \bar{L} \times (1+T)$; and none ($\bot$) otherwise.  
%
The purpose of $T$ is primarily to deal with short-term variations (jitter) in the observed performance of nodes,
reducing the alternation of "fast" and "slow" classifications for a same node. 
It achieves that by creating a neutral zone around the system average where proposal speed is considered neither "fast" nor "slow". The lower the $T$ the more non-$\bot$ decisions are achieved, at the cost of increasing sensitivity to transient latencies. The higher the $T$ the more resistant it is, but at the cost of less non-$\bot$ decisions. %$T$ must be a value $\geq0$, with an ideal minimum $T$ approximating how much consensus can vary due to noise in percents  ...???? 

\paragraph{Denote $c_h \in \{ fast, \bot, slow \}$} the above classification for consensus height $h$.

\subsection{M.2 - From local to global classifications}

Now, the challenge is to derive, from these local classifications, meaningful global decisions with respect to which nodes are considered fast, slow, or none.  
%
Our algorithm is based on three steps: \textit{vote}, \textit{choice}, and \textit{write}.   We propose here that these steps are carried out with the underlying consensus algorithm messages and steps, having an implicit trigger associated with it.  Our discussion uses Tendermint, with the following behavior at each node:

\begin{itemize}
\item \textit{Vote:} upon start of height $h+1$, send $\langle vote,h,c_h \rangle$ to all other nodes.  $c_h$ is the classification already discussed regarding consensus instance at height $h$.   
%
This vote is sent along with the first message for height $h+1$ to all other nodes  (e.g. Tendermint's prevote).  
Voting lasts until the decision for height $h+1$.  
The reason for using the decision as a limit is twofold. First, the synchronization eases the integration of the mechanism and reuse of messages. %, ensuring scores are updated in a timely manner. 
And second, it allows to gather additional votes beyond the quorum needed to advance the underlying consensus phase.

\item \textit{Choice:} upon start of height $h+2$, if a quorum of $2f+1$ \textit{vote} messages for a same non-$\bot$ $value$ is reached,  send a signed message $\langle choose,h,value \rangle_{sig}$ to all other nodes.   Otherwise, send a $\langle choose,h,\bot \rangle$ non-signed message. 
As with the voting step, the $choose$ message is sent along the first message for the next height ($h+2$)
and its reception lasts until the decision for that height.    


\item \textit{Write:}  Upon start of height $h+3$, nodes gather the chosen values and decide the value to write, called $decided_h$.   
%
If a node has received a quorum of $2f+1$ 
\textit{choose} messages for a same non-$\bot$ value,
the node decides that value. 
Otherwise decides the default value of $\bot$.

The proposer of height $h+3$ $write$s with the proposal: the $decided_h$ value; if $decided_h \neq \bot$, the quorum of signed choices; and the updated scores of all nodes.  Scores are discussed in Section \ref{subsec:scores}.  If $decided_h = \bot$, there is no need to record the quorum of signed choices.  
% COMPLEMENTAR ISSO ?
If a dishonest proposer node may omit, not writing a non-$\bot$ value, an honest next proposer may write the value with the quorum of signed choices justifying it.
Each proposer node thus keeps a window of $decided_h$ with their respective set of signed choices justifying the value.

Any node, upon receiving a proposal which can be accepted by the underlying consensus algorithm, adds the following: if $decided_h = \bot$, accepts it;  if $decided_h \neq \bot$, accepts if signatures are successfully verified. If not accepted, evidence of dishonest behavior is generated.

%A node only votes for the proposal (e.g. Tendermint's block) if the decided value and the updated scores in the proposal match its own values. This is necessary to prevent Byzantine proposers from simply ignoring the last vote. 

\end{itemize}

Note that all steps can happen simultaneously. During a consensus for height $h$, there can be an ongoing $vote$ step for height $h-1$, a $choice$ step for height $h-2$, and a $write$ step for $h-3$. %This only implies in two extra values (vote and choice) in each node's first message of that height. 


\subsection{M.3 - Continuous scoring and proposer prioritization}
\label{subsec:scores}

So far we have discussed how each consensus instance $h$, and it's associated proposer $n$, is associated with a classification $decided_h$ in $\{fast, \bot, slow \}$ known by all nodes.
%
Now we address how this classification can be used to steer proposer selection.   There are two main aspects to consider: 
(i) assuming a \textit{nextProposer} function using the $stake_n$ of each node $n \in \Pi$ to decide the proposer sequence, how can we extend the function to prioritize faster nodes;
(ii) considering that the relative performance of nodes may vary over time, due to several external factors, how the prioritization adapts to the changing relative performance of proposers.

\paragraph{We define $weight_n$ for each node $n \in \Pi$,} as: 
\begin{equation} 
\label{eq.weight}
weight_n = stake_n \times [1 + ( C \times nomalizedScore_n) ]
\end{equation}

Instead of inputting $stake$, the \textit{nextProposer} function is used with $weight$.   $C$ and $nomalizedScore_n$ are discussed next.


\paragraph{We define the input parameter $C$, the contribution of score,} to denote how much each node's score modifies its stake. This ensures that, in systems where different stake values are possible, the score remains proportional to those different stakes. 
$C=0$ disables the score algorithm and values $\geq1$ would allow nodes to be removed from the proposer pool. $0 < C < 1$ keeps all nodes in the choice and allows to differentiate them by performance.  The higher the $C$ the more performance is valued to choose the next proposer.  
%
For instance, $C=0.3$ would limit the $weight$ to $\pm30\%$  of each node's stake, leading to a $weight_n$ in the interval of $[0.7*stake_n, 1.3*stake_n]$. 

\paragraph{We now define $score_n$, $\alpha$, $\beta$ and $M$.}

We call $decided_n$ a value in $\{ -1, 0, 1 \}$ as $decide_h$ is respectively $\{ slow, \bot, fast \}$. 
A $score_n$ of each node $n \in \Pi$ is initialized to 0.
Each time a new $decided_n$ value is available, $score_n$ is updated as: 
\begin{equation} 
\label{eq.score}
score_n' = trunc_{[-M..M]}(\beta*score_n + \alpha*decided_n )
\end{equation}
%
\paragraph{$\beta$ denotes how fast older decisions decay their impact on score.}
$\beta$ is a value in $0..1$ 
%and serves to decrease the value of older decisions in favor of newer ones.   This is one measure to
and serves to adapt the prioritization of nodes to current observed performance.   It adjusts how quickly the score returns to zero with only $\bot$ decisions.

\paragraph{$\alpha$ denotes how new decisions impact the score.} 
$\alpha$ is a positive value, being a step in the modification of score. 
A higher $\alpha$ will move score faster, making score more sensitive to performance oscillations.   
If there are many non-$\bot$ decisions we can afford a small step, but if the number of non-$\bot$ decisions is small we are forced to have a larger $\alpha$. 
%
With these definitions, for any $0 \leq \beta < 1$  score is within the interval $\pm \frac{\alpha}{1-\beta}$.\footnote{Derivation of this expression omitted due to size restrictions.} 

\paragraph{$M$ is a variable which informs the limits of $score$,} which is restricted to the interval $[-M,M]$.   If $M = \frac{\alpha}{1-\beta}$ it allows the full variation of scores (truncation does not take place). 

%$M$ must be a value in the interval $[0, \frac{\alpha}{1-\beta}]$, although it should ideally be equal to that limit to allow for the full resolution of scores. 



\paragraph{We now define $nomalizedScore_n$}, resulting in a value in $-1..1$,  to be used in Equation \ref{eq.weight}, as: 
\begin{equation} 
\label{eq.normalizedscore}
nomalizedScore_n =  \frac{score_n}{M} 
\end{equation}
 



%\vspace{4mm}
As mentioned, the $nextProposer$ function needs to be deterministic in a decentralized environment. Therefore,  variables \(\alpha\), \(\beta\), $M$, $C$, and $T$ must be consistent across nodes. 
%
If the consensus supports self-configuration it is also possible to self-configure the global variables according to a pre-established set of rules, e.g. increasing the $T$ if there's too much jitter and decreasing it if there's too few decisions. For systems with self-configuration, modified variables must be stored (e.g. inside a blockchain's block) so that new nodes can update their variables to match the system. 


\subsection{Correctness and Byzantine Fault Tolerance}

The proposed mechanisms should ensure the properties described in the introduction of Section \ref{sec:algorithm}.
Each one discussed next.
% \begin{description} 
%     \item[P.1] the underlying consensus algorithm properties are preserved;     
%     \item[P.2] disregarding stake, or with a uniform stake distribution, proposer nodes with faster (slower) decisions become proposers proportionally more (less) often;
%     \item[P.3] changes in a node's relative performance over time lead to respective adaptation of its frequency as proposer;  
%     \item[P.4] every node eventually becomes proposer;  
%     \item[P.5] dishonest nodes do not impose an arbitrary sequence of proposers;
%     \item[P.6] dishonest nodes do not hinder the differentiation of faster/slower nodes. 
% \end{description}

\paragraph{Argument for P.1:} 
Regarding the properties of the underlying consensus algorithm, notice that the proposed mechanisms do not modify the algorithm.   Messages and the value proposed are conservatively extended with votes and scores,  these being processed by mechanisms separated from the underlying algorithm.  The proposed mechanisms also do not prevent the progress of the underlying algorithm: whenever 
phase or height change may occur, the proposed mechanisms proceed with the values computed so far.  

The only aspect impacted is the choice of next proposer.
The structure of the function to choose the next proposer is preserved, however being parameterized with weighted stakes according to respective nodes performance, instead of stakes only. As the weighted scheme, by definition, does not nullify any stake, every node is eventually chosen as proposer. 

\paragraph{Argument for P.2:} 
As discussed in Section \ref{subsec:scores} the $score_n$ of a node $n$ increases each time it is voted $fast$, leading to a multiplication of $stake_n$ by a factor $> 1$, that then results in proportionally augmented $weight_n$.   The weights of all nodes then substitute their original stakes in the $nextProposer$ function, leading to a more frequent choice of nodes according to their weight. 

\paragraph{Argument for P.3:} 
The proposed mechanisms continuously sample consensus latencies, adjusting scores.
The $\beta$ parameter serves to decay the importance of old samples in the current score while $\alpha$ denotes the impact of new decisions.  With this, scores adapt to changing performance observations, leading to the desired impact in the $nextProposer$ choice.

\paragraph{Argument for P.4:} 
This property derives from assumption A.1 and A.2 on the $nextProposer$ function, implying that nodes are chosen as proposers proportionally to their stake, and the fact that the scoring mechanism does not annulate any stake, but differentiates them according to the node's performance.  

\paragraph{Argument for P.5:} 
A non-$\bot$ value is written if $2f+1$ nodes choose that value.  Among those, at least $f+1$, a majority, is comprised by honest nodes, by definition.  This means that dishonest nodes cannot impose a non-$\bot$ value not supported by a majority of honest ones.


\paragraph{Argument for P.6:} 
Since there are at least $2f+1$ honest nodes, if they choose the same non-$\bot$ value, 
then this value is eventually written by an honest proposer irrespective of dishonest nodes choices
or dishonest proposers omitting the value.

\paragraph{Impact of dishonest behavior while voting:}

Properties P.5 and P.6 leave room for dishonest nodes impacting the decision if
the number of honest nodes choices $hc$ for a same non-$\bot$ value is $f+1 \leq hc < 2f+1$.

However, notice that dishonest nodes choices are limited: 
%In the \textit{vote} step, each node awaits for a quorum of other nodes. 
A node only chooses a non-$\bot$ value in the \textit{choice} step if it received a quorum of that value in the \textit{vote} step. 
As any two quorums have an overlap of at least one honest node, it is impossible for there to be two conflicting quorums. At the end of this step each node will either have a quorum of the same non-$\bot$ value, or default to $\bot$ if a quorum has not been received.  In the first case, a dishonest node sending a non-$\bot$ choice contradicting the quorum value would immediately generate evidence of dishonest behavior.
%This reduces the next step into a binary decision. 
%
This means that colluding dishonest nodes may either add choices to a majority of honest nodes 
to accept their non-$\bot$ value, or not, keeping the $\bot$ decision.      
Keeping the $\bot$ decision does not harm performance,  while accepting the non-$\bot$ value
would enhance the performance observed by a majority of honest nodes.

\paragraph{Impact of dishonest behavior while proposing:}
In the role of proposer, a written non-$\bot$ value is accompanied by the signed choices justifying the quorum.
Any modification of the agreed value would generate clear evidence of misbehavior by the proposer as
different signed choices from a same dishonest node would be available at honest nodes.

A dishonest proposer may write $\bot$ (which does not need the accompanying signed choices), 
omitting a quorum of a non-$\bot$ value.   However, any honest next proposer can write the omitted 
non-$\bot$ value and the respective signed choices, proving the value.   In practice, the omission 
by the dishonest node leads just to postponing the adoption of the value.



